\documentclass[11pt,reqno]{amsart}

\usepackage{amsmath,amssymb,amsthm}
\usepackage[margin=1.15in]{geometry}
\usepackage{microtype}
\usepackage{booktabs}
\usepackage[colorlinks=true,linkcolor=blue,citecolor=blue,urlcolor=blue]{hyperref}

\theoremstyle{plain}
\newtheorem{lemma}{Lemma}[section]
\newtheorem{proposition}[lemma]{Proposition}
\newtheorem{conjecture}[lemma]{Conjecture}
\theoremstyle{definition}
\newtheorem{remark}[lemma]{Remark}
\newtheorem{question}[lemma]{Question}
\newtheorem{step}{Step}

\newcommand{\R}{\mathbb{R}}
\newcommand{\Z}{\mathbb{Z}}
\newcommand{\F}{\mathbb{F}}
\newcommand{\Inat}{I^{\natural}}
\newcommand{\HF}{\mathit{HF}}
\newcommand{\CF}{\mathit{CF}}
\newcommand{\rk}{\operatorname{rank}}
\newcommand{\sgn}{\operatorname{sgn}}
\newcommand{\pill}{P}
\newcommand{\Rnat}{\mathcal{R}^{\natural}}
\newcommand{\Rp}{\mathcal{R}}
\newcommand{\iss}{\looparrowright}
\newcommand{\bc}{b}
\newcommand{\Qh}[1]{\widehat{Q_{#1}}}

\begin{document}

\title[Finding the pillowcase bounding cochains]{Finding the pillowcase bounding cochains:\\
the search behind the $P(-2,3,q)$ computations,\\ in complete detail}

\author{Bernd J. Wuebben}
\date{July 30, 2026}

\begin{abstract}
A correspondence note, written for Zuyi Zhang, expanding \S\S4--5 of the pretzel paper
\cite{P2} to full operational detail. Part I (\S\S\ref{sec:problem}--\ref{sec:yours}) is the
answer to ``the argument in the paper is too brief to follow'': the complete strategy by which
the bounding cochains of $P(-2,3,q)$ were found, including the parts that are usually left in
the source code --- the lift bookkeeping, the pruning discipline, the traps that produced wrong
answers along the way, and the acceptance criteria. It is written so that the same search can
be run on other pillowcase pairs, in particular on Kai Smith's examples \cite{Smi} where the
naive homology differs from $\Inat$. Part II (\S\ref{sec:analytic}) collects what I can currently say about
the analytic question we are both stuck on --- identifying the immersed polygons that carry the
higher operations with objects on the quilted side --- including two proposed reductions that
would sidestep the polygon-lifting problem for the deformed differential, and the smallest
concrete test cases I can offer. Nothing in Part I is new relative to \cite{P2} and its
ancillary code; every number quoted here is re-derived by that code.
\end{abstract}

\maketitle

\tableofcontents

%=====================================================================
\part{The search, step by step}
%=====================================================================

\section{The problem, and the one precondition}\label{sec:problem}

\subsection{The problem}
Fix a Conway-sum decomposition $K=(B^3,T_1)\cup_{(S^2,4)}(B^3,T_2)$ and the two immersed
Lagrangians in the pillowcase $\pill=(\R/2\pi\Z)^2/\iota$, $\iota(\gamma,\theta)=(-\gamma,-\theta)$:
in the pretzel family, $\mathrm{red}=\Rnat(\Qh{-1/2})$ (the earring tangle, whose own cochain
vanishes by Smith \cite[Lemma~5.6]{Smi}) and $\mathrm{blue}=\Rp_t(Q_{1/3}+Q_{1/q})$; here
$\Rp(T)$ is the traceless-character image of a tangle in the pillowcase, $\Rp_t$ its
holonomy-perturbed version, and $\Rnat$ the earring-added version ($\Qh{r}$ likewise denotes
$Q_r$ with the earring added), in the notation of \cite{P2}. The naive Lagrangian--Floer homology of the pair generally differs from
$\rk\Inat(K)$, Kronheimer--Mrowka's reduced singular instanton knot homology, and the
Cazassus--Herald--Kirk--Kotelskiy conjecture \cite{CHKK} says a bounding cochain
$\bc\in\CF(\Rp(T_2),\Rp(T_2))$ repairs it. The computational problem is:
\begin{quote}
\emph{find the minimal $\F_2$-combination $\bc$ of self-intersections of blue such that
$\bc$ satisfies the Maurer--Cartan equation to the computed orders (triangles and
quadrilaterals, \S\ref{sec:deformed}) and
$\rk\HF(\partial_\bc)=\rk\Inat(K)$.}
\end{quote}

\subsection{The precondition that makes the search well posed}\label{ss:precondition}
The single most important structural point --- the one I would fix first for any new example ---
is that \emph{the target rank must be known independently, before any pillowcase computation
happens}. Two reasons.

First, the Maurer--Cartan equation alone does not pin $\bc$: $\bc=0$ always satisfies it, and
(as $q=11$ shows, \S\ref{ss:q11}) even the rank condition can admit dozens of solutions. The
search needs an external answer to converge to.

Second, without a target there is no stopping criterion. The search enumerates supports of
increasing size (\S\ref{sec:search}); ``no solution of size $\le r$'' is only meaningful
relative to a rank one is trying to hit.

For the pretzel family the target is Theorem~1.1 of \cite{P2}:
$\rk\Inat(P(-2,3,q))=q+2$ for all odd $q\ge3$, proved unconditionally by squeezing
\[
   \ell(K)=\textstyle\sum_i|a_i(\Delta_K)|\;\le\;\rk\Inat(K)\;\le\;\dim\mathit{Kh}_r(K)
\]
between the Alexander lower bound of Kronheimer--Mrowka/Lim and the Khovanov upper bound of
Kronheimer--Mrowka, with both ends computed in closed form (a Conway-skein Chebyshev recursion
for $\Delta$; Manion's theorem for $\mathit{Kh}_r$). The squeeze is exportable: for any knot
where $\sum_i|a_i(\Delta)|=\dim\mathit{Kh}_r$ --- which includes many small knots and, by
Manion's closed forms, all $3$-strand pretzels $P(-2,3,q)$ --- the instanton rank comes for
free. For Kai's examples this is the first thing I would check; if the squeeze does not close,
some independent determination of $\rk\Inat$ is needed before the cochain search can be
trusted to mean anything.

\subsection{The direction of the correction}
Over $\F_2$, homology rank moves in steps of $2$: a deformation either \emph{cancels} an
existing entry of the differential (rank $\HF$ goes up by $2$) or \emph{creates} a new
independent one (down by $2$). So the sign of $\text{naive}-\Inat$ dictates the mechanism
before the search starts: deficit $\Rightarrow$ the added polygon terms must land exactly on an
occupied entry of $\partial_0$ and remove it; excess $\Rightarrow$ they must land on an empty
entry. In the pretzel family the sign is governed by the determinant:
$\text{naive}-\rk\Inat=2\,\sgn(\det K-3)$ at all five computed members, with $\det=|q-6|$.
The deficit regime $\det<3$ is exactly the regime where the double branched cover is an
integral homology sphere.

\section{The frame: the double cover, and the two bookkeeping rules}\label{sec:frame}

All computations happen on the torus double cover $T^2\to\pill=T^2/\iota$, with curves stored
as $\iota$-invariant closed polylines (a few hundred vertices each; blue at $q=5$ has $460$
edges after adaptive refinement). The pillowcase itself is used only for reporting
coordinates. Two bookkeeping rules make the cover counts equal the pillowcase counts, and both
were hard-won.

\subsection{The $C_+$ lift rule}\label{ss:cplus}
Red lifts to two $\iota$-related circles on $T^2$; fix one of them, $C_+$, once and for all.
A pillowcase polygon lifts to exactly two $\iota$-related $T^2$ polygons, and exactly one of
the two has its red arc on $C_+$. Therefore: \emph{count on $T^2$, keeping only polygons whose
red arc lies on $C_+$, and the result is the pillowcase count --- no division by two, no
orbifold correction}. This is Lemma-level, but it is the reason the whole computation can be run
on the cover, where the winding-number machinery of \S\ref{sec:naive} is available.

A pillowcase self-intersection of blue, by contrast, has \emph{two} $T^2$ preimages, and a
polygon may close through either; contributions of both preimages are accumulated. (Missing
the second preimage silently halves some triangle counts; the code carries both throughout.)

\subsection{The folding trap: locality lives in edge index, not in distance}\label{ss:folding}
(The fold here is the quotient $T^2\to\pill$ --- no relation to the quilt folding trick of
Part~II.) An immersed polygon among our curves is a small disk, so its vertices --- the two
generators and the $\bc$-crossings (call a self-intersection of blue a \emph{$\bc$-crossing}
when it is a candidate summand of the support of $\bc$) --- sit in one geometric cluster, and
enumeration is pruned accordingly. The trap: \emph{some generators fold under $T^2\to\pill$}. At $q=5$, four of the
nine generators lie near $\theta\approx1.5$ on the cover but at $\theta\approx4.7$ in $\pill$;
points close in $\pill$ can be far apart on the cover and vice versa. A pillowcase-metric
pruning is therefore wrong on both sides, and in an intermediate version of this work it
silently discarded a genuine quadrilateral --- the $q=5$ cochain quadrilateral itself --- and
produced ``no cochain exists'' as the answer. The correct locality test is the
\emph{circular edge-index window}: the blue boundary of a small immersed disk traverses a
contiguous band of blue edges, so the anchor edges of all vertices (each generator's blue edge
index, each crossing's two branch indices) must fit inside a circular window on the blue edge
cycle. The window used in every reported computation is $230$ edge indices, deliberately loose --- pruning is a
tractability device, never a correctness device, and the triangle enumeration
(\S\ref{ss:tables}) runs unpruned as a check on exactly this class of error.

\section{The naive complex}\label{sec:naive}

\subsection{Generators and the bigon predicate}
Generators are the transverse intersections $\mathrm{red}\cap\mathrm{blue}$ (fiber products on
the cover, in Zhang's sense --- distinct preimages are distinct generators). A candidate bigon
from $x$ to $y$ is a closed loop $\Gamma=\alpha\cdot\bar\beta$, $\alpha\subset$ red from $x$
to $y$, $\beta\subset$ blue from $y$ to $x$. On the cover, let $w$ be its winding function.
Then $\Gamma$ bounds an immersed bigon iff
\begin{itemize}
\item[(a)] $[\Gamma]=0\in H_1(T^2)$;
\item[(b)] $w=0$ at (the preimages of) the four pillowcase corners;
\item[(c)] $w\ge0$ on every complementary face, and $w>0$ somewhere;
\item[(d)] the corners at $x$ and $y$ are convex: a single embedded sector of angle $<\pi$,
across which $w$ jumps by exactly $+1$.
\end{itemize}
The differential is $\partial x=\sum_y(\#\{\text{bigons }x\to y\}\bmod2)\,y$ and
$\rk\HF=\#\text{gens}-2\rk_{\F_2}\partial$.

(For embedded curves this winding-number criterion is de~Silva--Robbin--Salamon's
combinatorial Floer homology; the immersed extension used here is that of \cite[\S4]{P2}.)

Implementation notes that matter in practice: $w$ is evaluated at
probe points offset from each edge midpoint on both sides, with the offset adapted to the
clearance from the other edges --- a fixed offset either steps over a thin face or lands on a
neighboring edge; near-parallel arcs are the norm here, not the exception (at $q=5$ the
interior blue strands cross red at angle $\approx0.026$); corner convexity uses an explicit
sector test rather than a sign heuristic, precisely because of those shallow angles.

\subsection{Validation before use}\label{ss:gates}
The programs are validated against everything the literature provides before any new number is
computed. At $q=5$ the reconstruction reproduces Smith's published output \cite{Smi} in every particular:
nine generators (one interior, two clusters of four), exactly two bigons, naive rank $5$, and
his worked example of the seam-circle resolution (the diagonal cut-and-paste that replaces each
seam circle of reducibles by embedded arcs with one new self-crossing). The polygon counter
(\S\ref{ss:tables}) reproduces
the one known higher-polygon mechanism --- the single $\mu^2$ triangle through the earring's
pinch self-crossing, on the unlink example Smith uses as his simplest test case --- and is
exercised on synthetic teardrop and lens configurations
where the right answer is forced. Only after these checks pass is anything new computed.

\subsection{The perturbation discipline}\label{ss:pert}
Both curve constructions carry perturbation parameters (the seam-circle resolution amplitude
for blue; the earring perturbation for red), and nothing is reported from a single
perturbation. The rules:
\begin{itemize}
\item Every rank and every cochain is recomputed at two independently chosen perturbations.
Generator \emph{counts} are not invariants and do vary ($19$ vs $23$ at $q=11$; $25$ vs $27$
at $q=13$); the homology ranks must agree, and do in every reported case.
\item A cochain is accepted only if its support recurs as \emph{the same geometric crossing}
at both perturbations, within the perturbation amplitude. The $q=7$ cochain, for instance,
recurs at $(0.042,5.405)$ and $(0.058,5.413)$.
\item Degenerate perturbations exist and must be screened for: one non-generic earring
perturbation at $q=7$ collapses the thirteen generators to nine. The screen is exactly the
two-perturbation comparison --- a collapse shows up as a generator-count discrepancy that
persists under refinement.
\end{itemize}

\section{The deformed differential as a table}\label{sec:deformed}

\subsection{The polygon tables}\label{ss:tables}
The bigon predicate generalizes verbatim to $(k{+}2)$-gons: the boundary loop now alternates
between red and blue arcs and switches branches at self-crossings of blue, the corner test of
\S\ref{sec:naive}(d) is applied at \emph{every} vertex, and conditions (a)--(c) are unchanged.
With $n$ generators and the $\bc$-crossings
ranged over the pillowcase self-crossings of blue ($40$ at $q=5$; $82$ at $q=7$; $225$ at
$q=11$), two tables are precomputed:
\begin{align*}
\mathrm{Tri}[S]_{xy}&=\#\{\text{immersed triangles }x\to y\text{ with }\bc\text{-vertex }S\}
\bmod 2\quad (n\times n\text{ over }\F_2),\\
\mathrm{Quad}[\{S,S'\}]_{xy}&=\#\{\text{quadrilaterals }x\to y\text{, }\bc\text{-vertices }
S,S'\text{, both cyclic orders}\}\bmod 2.
\end{align*}
Each entry accumulates over both $T^2$ preimages of each crossing (\S\ref{ss:cplus}). The
triangle sweep is exhaustive (no pruning); the quadrilateral sweep is pruned by the circular
edge-index window of \S\ref{ss:folding} only. For any support $\bc=\{S_i\}$,
\[
   \partial_\bc \;=\; \partial_0
   \;\oplus\;\bigoplus_{S\in\bc}\mathrm{Tri}[S]
   \;\oplus\;\bigoplus_{\{S,S'\}\subseteq\bc}\mathrm{Quad}[\{S,S'\}] ,
\]
evaluated by table lookup, so the search of \S\ref{sec:search} costs one $\F_2$ rank
computation per candidate. Higher orders were not needed at any computed member: at $q=5$ the
question ``are $\mu^4$ pentagon terms required?'' was answered by construction --- the
tri$+$quad tables already contain a rank-target solution, and a pentagon enumeration built as
a check found nothing that changes it. If a new example fails at tri$+$quad order, the same
architecture extends: the pentagon table is indexed by crossing triples.

\subsection{The Maurer--Cartan data}\label{ss:mc}
The same counter, run with \emph{all} vertices at self-crossings of blue, produces the
Maurer--Cartan data: monogons ($\mu^0$), self-bigons ($\mu^1$ on $\CF(\mathrm{blue},
\mathrm{blue})$, stored as a crossing-by-crossing matrix), self-triangles ($\mu^2$). Two
structural facts:
\begin{itemize}
\item $\mu^0=0$ always, for the standard degree reason (Proposition~4.2 of \cite{P2}): a
teardrop would sit in degree $2$ of the self-Floer complex, but a transverse self-intersection
of a curve in a surface contributes generators only in degrees $0$ and $1$ (Akaho--Joyce).
This is what makes the Maurer--Cartan equation a genuine equation. The counter nevertheless
computes monogon counts directly ($0$ everywhere, and it does detect a genuine teardrop on a
synthetic test curve).
\item $\mu^1$ is \emph{not} always zero: at $q=5$ it vanishes (and all
window-pruned self-triangle triples with it, so Maurer--Cartan is vacuous on small support),
but at $q=7$ and $q=11$ blue carries two self-bigons each, and $\mu^1$-closedness of the
support becomes a real constraint that kills candidates in the search.
\end{itemize}

\section{The search}\label{sec:search}

\subsection{The search procedure}
\begin{step}[Active set]
Restrict to crossings carrying at least one generator-touching polygon (an entry in some
$\mathrm{Tri}$ or $\mathrm{Quad}$ table). A crossing outside this set cannot change
$\partial_\bc$; adding it to a support only risks new Maurer--Cartan violations. Minimal
solutions therefore live in the active set. (At $q=11$: $108$ of $225$ crossings are active
through triangles.)
\end{step}
\begin{step}[Increasing support size]
For $r=1,2,3,\dots$: enumerate all $r$-subsets of the active set, evaluate
$\rk_{\F_2}\partial_\bc$ by table lookup, and keep the subsets hitting the target rank
$(n-\rk\Inat)/2$.
\end{step}
\begin{step}[Maurer--Cartan filter]
For each rank solution check, in order: $\mu^0=0$ at each support crossing; then, at
\emph{every} crossing $s$ of blue, the component of $\mu^1(\bc)+\mu^2(\bc,\bc)$ at $s$
vanishes --- the $\mu^1$ part from the self-bigon matrix, the $\mu^2$ part by counting
self-triangles with two vertices in the support and the third at $s$, over all orderings and
both preimages.
\end{step}
\begin{step}[Stop at first success; minimality by exhaustion]
The loop stops at the first size $r$ admitting Maurer--Cartan-valid rank solutions; minimality
is then a statement about exhaustion, not heuristics --- every smaller support was tested and
failed, and the failures are reproducible.
\end{step}
\begin{step}[Perturbation stability]
Re-run everything at the second perturbation (\S\ref{ss:pert}); accept only supports that
recur as the same geometric crossings.
\end{step}

\subsection{Near-misses are informative}
The search log is worth reading, not just the answer. At $q=5$ the geometrically most tempting
candidate --- the pair of crossings nearest the surviving bigon --- reaches the right entry but
drags three triangle terms along and lands at rank $3$, not $1$. That failure is what forced
the recognition that cancellation is an \emph{exact-incidence} condition: the added polygon
terms must land on the occupied entry and nowhere else. It is also why a search that stops at
triangles is wrong at $q=5$: no triangle through any of the forty crossings reaches either
bigon of $\partial_0$ --- the correction is forced one polygon order up, to the quadrilateral.
A triangles-only search reports ``no cochain exists''; this is the single most likely failure
mode I would expect on a new example.

\section{The three runs, with all the numbers}\label{sec:runs}

\subsection{$q=5$ (deficit; unique; quadrilateral)}\label{ss:q5}
Nine generators; bigons at entries $(x_1,y_0)$ and $(x_4,y_6)$; naive rank $5$; target
$\rk\partial_\bc=1$. Size-$1$ search: no single crossing achieves rank $1$. Size-$2$: exactly
one pair does, $\bc=s_A+s_B$ with $s_A$ at $(\gamma,\theta)\approx(0.028,1.272)$ and $s_B$ at
$(3.057,4.981)$,
one on each pillowcase seam, exchanged by the symmetry
$(\gamma,\theta)\mapsto(\pi-\gamma,-\theta)$. Mechanism: a single immersed quadrilateral with
generator vertices $x_4,y_6$ and $\bc$-vertices $s_A,s_B$; neither crossing carries any other
generator-touching polygon, so $\partial_\bc=\partial_0+E_{x_4y_6}$, leaving the single entry
$(x_1,y_0)$: $\rk\HF=9-2=7=\rk\Inat$. Maurer--Cartan is vacuous here ($\mu^1=0$, $\mu^2=0$ on
all tested triples).

\subsection{$q=7$ (deficit; unique; triangle)}\label{ss:q7}
Thirteen generators; three bigons; naive rank $7$; target $\rk\partial_\bc=2$. Here a single
crossing suffices: $s$ at $(\gamma,\theta)\approx(0.05,5.41)$ carries an immersed triangle whose generator pair
coincides with one bigon of $\partial_0$, carries no other generator-touching polygon, and is
$\mu^1$-closed (blue has two self-bigons at $q=7$, so this is a real constraint). The
size-$1$ search, Maurer--Cartan filtered, returns exactly this crossing, at both
perturbations. $\rk\HF=13-4=9=\rk\Inat$. The polygon order of the correction is thus not
constant even within the deficit regime: quadrilateral at $q=5$, triangle at $q=7$.

\subsection{$q=11$ (excess; abundant; triangles)}\label{ss:q11}
Nineteen generators at the reported perturbation; naive rank $15>13=\rk\Inat$; the correction
must \emph{create} a differential, raising $\rk\partial_\bc$ from $2$ to $3$. Creation is
combinatorially generic: any triangle through a single crossing whose generator pair is not
already occupied adds an entry. Of $225$ crossings, $108$ carry triangles; after the rank and
Maurer--Cartan filters, \emph{fifty-five} single-crossing bounding cochains remain, each with
$\rk\HF=19-6=13$. Fifty of the fifty-five created differentials flow into a single common
target generator. The asymmetry against \S\ref{ss:q5}--\ref{ss:q7} is structural:
cancellation is an exact-incidence condition and pins the cochain uniquely; creation is an
avoidance condition and manifestly does not. Rank-matching therefore cannot select \emph{the}
cochain of the tangle at excess members --- whatever universal rule attaches the cochain (in
the sense of your Remark~6.2, on which more in \S\ref{ss:universal}) must use more than the
rank.

\section{Running it on your and Kai's examples}\label{sec:yours}

The search needs three inputs: (i) the two curves, either as closed polylines on the $T^2$
cover or as tangle data from which the representation-variety construction can rebuild them
(for Conway sums of rational tangles the code does this from the fractions alone); (ii) a
validation anchor --- any independently known count for one member (a generator count, a bigon
count, a published rank) to anchor the reconstruction; (iii) the target rank
(\S\ref{ss:precondition}). Cost is minutes per member on a laptop, pure Python standard
library, the quadrilateral sweep dominating.

Everything is public at
\url{https://github.com/bwuebben/pillowcase-bounding-cochain}
(and ships as ancillary files with the arXiv postings \cite{P1,P2}); Appendix~\ref{app:files} maps each
step above to its module. If any of Kai's homology-discrepancy examples comes with an
independently known instanton rank, I would genuinely like to run the search on it --- and if
the naive/instanton discrepancy there exceeds one differential, the outcome (minimal support
size, uniqueness or abundance) would be the first data point beyond $|{\det}-3|$-governed
behavior of the pretzel family.

%=====================================================================
\part{The analytic question}
%=====================================================================

\section{On polygons, quilts, and the folding trick}\label{sec:analytic}

This part responds to your question --- how to identify the polygons carrying the higher
operations with objects on the quilted side, given that the folding trick is built for strips
--- with the caveat stated plainly: I do not have a proof either. What I can offer is (i) a
separation of the problem into independent pieces, (ii) two proposed reductions that would
remove the piece my computations actually need without ever lifting a polygon, (iii) a precise
map of where your strip machinery would need to change if one lifts polygons directly, and
(iv) concrete test cases.

\subsection{Separating the two gaps}\label{ss:twogaps}
Two logically independent things are unproved, and they fail differently.
\begin{itemize}
\item[(G1)] \emph{Faithfulness of the combinatorial model} (Hypothesis~4.1 of \cite{P2}): the
winding-number counts of \S\ref{ss:tables} compute the $\mu^k$ of the immersed Fukaya algebra
of $\pill^*$, the pillowcase with its four orbifold points removed. This is a statement about one surface; no correspondence, no quilt appears in
it. Your uniqueness theorem is its $k=1$ case on closed surfaces.
\item[(G2)] \emph{The instanton comparison} (Conjecture~1.1 of \cite{CHKK}): the $\bc$-deformed Floer
homology is $\Inat$. This is where quilts, strip shrinking, and figure-eight bubbling live ---
and where your Theorem~5.1 (strips upstairs from strips downstairs) and the Bottman--Wehrheim
program sit.
\end{itemize}
Everything computed in Part I is a statement inside (G1)'s model and is agnostic about (G2).
The polygon-lifting question you raised is (G2)-flavored, but --- and this is the point of the
two reductions below --- the deformed \emph{differential} may not need polygon lifting at all.

One more separation, within (G1) itself, because your two papers divide exactly along it. A
combinatorial computation of an operation needs three legs: \emph{existence} (every
combinatorial polygon has a holomorphic representative --- your construction paper),
\emph{uniqueness} (it has only one --- your uniqueness paper, within a smooth isotopy class
with prescribed boundary image), and \emph{completeness} (there are no holomorphic solutions
outside the combinatorial list). The third leg is the one CHKK flag themselves at the end of
\cite[\S10.4]{CHKK} --- ``there may be quilts or polygons which we have overlooked'' --- and it is the leg
their conjecture actually turns on, since ``this is the only differential'' is a statement
about the whole moduli space, not about one homotopy class. Any use I make of your theorems
below silently needs completeness too; I flag it here once rather than at every step.

\subsection{Reduction 1: surgery instead of higher polygons}\label{ss:surgery}
The deformed differential $\partial_\bc$ is expected to be computed by an ordinary bigon count
after Lagrangian surgery. Concretely: let $L$ be an immersed curve, $s$ a transverse
self-intersection, and $L_s$ the smoothing of $s$ in the sector distinguished by the degree
convention on $\CF(L,L)$ (of the two smoothings of the crossing, the one rounding the sectors
a polygon corner at $s$ occupies). The expectation, standard in spirit since
Fukaya--Oh--Ohta--Ono's disk-trading analysis and proved in dimension $\ge2$ by
Palmer--Woodward \cite{PW}, is
\[
   (L,\bc=s)\ \simeq\ (L_s,0)
   \qquad\text{in the Fukaya category,}
\]
so that for any test object $R$: $\HF((R,0),(L,s))\cong\HF(R,L_s)$. On the combinatorial side
this has a precise, checkable shadow:
\begin{lemma}[combinatorial surgery lemma --- proposed]\label{lem:surgery}
For transversely immersed curves $R,L$ on a surface and a self-crossing $s$ of $L$, the
winding-number bigon counts of the pair $(R,L_s)$ equal the counts of the combinatorial
$\partial_{\bc=s}$ of \S\ref{ss:tables}: bigons of $(R,L)$, plus, for each $k\ge1$, the
$(k{+}2)$-gons of $(R,L)$ all of whose $\bc$-vertices ($k$ of them) lie at $s$ (triangles
through $s$ once, quadrilaterals through $s$ twice, and so on).
\end{lemma}
The geometric content is local and elementary: smoothing the crossing rounds the corner, so a
triangle with corner at $s$ becomes a bigon of $L_s$ passing once through the surgery site,
a quadrilateral with corners at two smoothed crossings becomes a bigon passing both sites,
and bigons avoiding the sites are untouched --- while a bigon of $\partial_0$ cancelled by the
deformation reappears \emph{paired} with a polygon-turned-bigon and dies mod $2$. I have not
written the general proof, but the lemma is \emph{verified entrywise at both deficit
members} by direct recomputation (Appendix~\ref{app:files}): smoothing blue
$\iota$-equivariantly at the support crossings --- both $T^2$ preimages, both sector
choices --- and recounting bigons against red reproduces $\partial_\bc$ \emph{entrywise}. At
$q=5$ the surgered pair's bigon matrix is exactly $\{(x_1,y_0)\}$ with $\HF=7$; at $q=7$ it is
exactly the two surviving entries with $\HF=9$. The check also makes the degree bookkeeping
visible: at $q=7$ one of the two smoothings of the crossing realizes the deformed complex and
the other reproduces the naive one --- the two smoothings are the two generators of
$\CF(L,L)$ at the crossing, behaving exactly as the sector convention predicts.

Why this matters for your question: \emph{if} Lemma~\ref{lem:surgery} holds and \emph{if} the
surgery equivalence holds for curves on surfaces, then the faithfulness of the deformed
differential --- the only thing my computations need from (G1) beyond your theorem --- follows
from the $k=1$ case alone. The chain is: combinatorial $\partial_\bc$ of $(R,L)$
$\;=\;$ combinatorial $\partial$ of $(R,L_s)$ [the lemma] $\;=\;$ analytic $\partial$ of
$(R,L_s)$ [your uniqueness theorem, applied to the surgered pair] $\;=\;$ analytic
$\partial_\bc$ of $(R,(L,s))$ [surgery equivalence]. No polygon is ever lifted; the higher
operations enter only through the surgered curve's bigons, which your theorem covers.

Two caveats, stated so you can attack them. First, Palmer--Woodward prove surgery invariance
in dimension $\ge2$; curves on surfaces are exactly the excluded edge case, so the analytic
input is genuinely open here --- though the same low dimension is what makes a direct
combinatorial proof plausible. Second, the sector bookkeeping (which smoothing corresponds to
which degree convention at $s$) is exactly the kind of detail that flips signs in odd ways;
over $\F_2$ it degrades to a two-case check, but the two cases must be matched to the two
generators of $\CF(L,L)$ at $s$ correctly.

\subsection{Reduction 2: the orbifold points, via the double cover}\label{ss:orbifold}
The second extension you would need --- past the four $\Z/2$ points of $\pill$ --- I believe is
the milder one, and \S\ref{ss:cplus} is the reason. The precise statement that would suffice:
\begin{quote}
\emph{Let polygons be constrained to a compact subsurface
$\pill_\rho\subset\pill^*$ obtained by deleting $\rho$-disks around the four orbifold points.
Then the immersed Floer theory of curves in $\pill_\rho$, with an $\iota$-invariant almost
complex structure pulled back from $T^2$, is computed on the free double cover
$T^2\setminus(4\text{ points})\to\pill_\rho$, with the $C_+$ normalization of
\S\ref{ss:cplus} converting cover counts to base counts.}
\end{quote}
Every relevant object in the computed examples clears the corners by a definite margin: all
self-crossings of every computed member lie at distance $\ge0.23$ from every corner (the
minima are $0.61$, $0.42$, $0.27$, $0.23$ for $q=5,7,11,13$), the cochain-carrying crossings
at distance $\ge0.27$, and the two deficit-member cochains at $\ge0.87$. So $\rho$ can be
taken macroscopic. What the reduction does \emph{not} handle by itself is a compactness
statement --- a limit of holomorphic disks escaping into the deleted disks. Here the concrete
geometry should help: near a corner the two curves are a pair of near-straight arcs through
the fold, and a disk limiting into the corner would have to carry area bounded below by the
corner clearance, giving an energy threshold below which escape is impossible. That is an
argument sketch, not a proof; but it is the reason I expect the closed-surface picture (i.e.\
your setting) to be the true home of the problem, with the orbifold a bounded-geometry
complication rather than a conceptual one.

\subsection{If one does lift polygons: where the strip enters, on both sides}\label{ss:polylift}
Reading your two papers with the polygon question in mind, I tried to list exactly where
``the domain is a strip'' is used essentially --- because that list is the specification for
any polygon version. The two halves break in different places.

\subsubsection*{The construction half}
In \cite{Z1}, as I read it:
\begin{enumerate}
\item The constant extension that seeds the implicit function theorem ($u_2'$, $u_1$, $u_1'$
constant in $t$ on $\R\times[0,\delta]$) needs the seam to be a full boundary line with a
product collar. A $(k{+}2)$-gon presented as a disk with boundary punctures has strip-like
ends, and between punctures each boundary arc has a product collar --- so the constant
extension survives \emph{locally along each arc}, and the question concentrates at the
punctures and corners.
\item The $\delta^{1/4}$ smallness of $\mathcal F_u(0)$, the $\delta$-uniform trace estimate,
and the integration by parts with boundary terms on disjoint horizontal lines all use the
product structure; on a punctured-disk domain the same estimates should localize to the collar
pieces, with the corner contributions entering through the strip-like ends' asymptotics.
\item The model operator at the ends ($\partial_s-J(x^\pm)\partial_t$ with Robbin--Salamon
Fredholm theory) generalizes to $k+2$ ends; what changes is bookkeeping, not analysis.
\item The index count does change: your Claim~5.9 uses index $1$ and the $\R$-translation to
get a bounded right inverse on a complement of the kernel; a rigid polygon has index $0$ and
no automorphism, which if anything simplifies that step.
\end{enumerate}
So my reading is that the analytic core of your \S5 is not strip-bound so much as
\emph{collar-bound}, and the genuinely new difficulty for polygons is (i) the corner
asymptotics at the $\bc$-vertices --- whose upstairs meaning is exactly the fiber-product
identification of your Proposition~2.28, extended from generators of $\CF(L_1,F\circ L_2)$ to
\emph{self}-intersections of the composed curve --- and (ii) the question I could not answer
from \S5 as written even in the strip case: the identification
$\mathcal X^Q\cong\mathcal X^Q_\delta$ across seam widths. The ratio of patch widths looks
like a conformal modulus of the two-patch quilt, invariant under simultaneous rescaling, so I
do not see how a width-$\delta$ quilt is identified with a width-$1$ quilt without deforming
the almost complex structure along the way --- am I misreading the intended argument? I ask
partly because the same point returns with force for polygons, where the seam position
acquires genuine moduli.

\subsubsection*{The uniqueness half}
In \cite{Z2} the wall is sharper, and it is conformal. Your proof runs on exactly three
boundary normalizations: the two strip ends, plus the point supplied by the
$\R$-translation normalization $u_2^t(0,2)=w_2^t(0,2)$; after uniformizing the
keyhole-excised domain and its image to the disk, three boundary points is precisely the
dimension of $\operatorname{Aut}(\mathbb D)$, and rigidity does the rest. A $(k{+}2)$-gon
domain has no translation action and carries $k-1$ real conformal moduli (cross-ratios of the
boundary punctures), so two competing holomorphic polygons with the same corners need not
even live on the same uniformized domain, and the three-point argument has no analogue. But I
do not think this kills the program --- I think it says the statement must be reformulated.
On a surface, the domain modulus of a rigid holomorphic polygon is not free: it is determined
by the image region, because the region (an immersed disk with corners) carries its own
conformal structure and the map is a uniformization of it. So the correct $k\ge2$ statement
is presumably not ``unique for each marked domain'' but:
\begin{quote}
\emph{each combinatorial polygon (immersed disk with convex corners at the prescribed
generators and crossings) is represented by a holomorphic $(k{+}2)$-gon for exactly one value
of the domain modulus, uniquely up to reparametrization} ---
\end{quote}
uniqueness with the modulus floating, one rigid solution per region. That formulation
restores the count $=\#\{\text{regions}\}$, which is what the combinatorial model uses, and
it seems susceptible to exactly your keyhole-plus-uniformization method, run with the domain
cross-ratios as unknowns rather than as givens. The quilted version acquires one more layer:
the seam pattern of a quilted polygon (which patch meets which arc, in which cyclic order) is
itself combinatorial data, and its degenerations are governed by the $2$-associahedra of
Bottman; Ma'u--Wehrheim--Woodward's $A_\infty$-functor machinery \cite{MWW} is where quilted
polygon domains are set up systematically, and may save you from reinventing the domain
theory.

\subsection{Perhaps polygons need not be lifted at all}\label{ss:broken}
There is a reading of CHKK \S10.4 under which your question --- \emph{what} in the quilts
corresponds to the downstairs polygons --- already has a proposed answer, and it is not a
quilted polygon. In their degeneration picture \cite[Figures 22--23]{CHKK}, the object upstairs is
a quilted \emph{bigon}; what happens under strip shrinking is that it limits onto a
\emph{triangle downstairs with a figure-eight bubble attached}, the bubble outputting the
cochain element. Read as a general mechanism, this says: the downstairs $\mu^2$-triangle
through a crossing $s$ is not the shadow of a quilted triangle --- it is a \emph{boundary
stratum of the moduli of quilted bigons}, the stratum where a figure-eight bubble forms with
output $s$. Under the Bottman--Wehrheim proposal (cochains from figure-eight bubbles), the
$\bc$-deformed downstairs differential would then be exactly the count of quilted-bigon ends:
the interior of the moduli contributes the naive bigons, the figure-eight boundary strata
contribute the $\mu^{k\ge2}$ corrections, and the deformed complex is the statement that the
two ends of a compact one-manifold cancel. If that is the right picture, the polygon-lifting
problem dissolves: one never lifts a triangle; one proves instead that the figure-eight
boundary strata of the bigon-quilt moduli are in bijection with the combinatorial
triangles-through-$s$. Your Example~3.6 mechanism in \cite{Z2} (the mean-value argument
forcing blow-up at the bifold, in your terminology) reads like the existence half of exactly that bijection, at
the first interesting crossing. I do not know how to prove the bijection in general --- but
it has a checkable shadow, namely Lemma~\ref{lem:surgery}, and the two computed cancelling
cochains of Part~I are the two smallest instances in which the bijection's count is forced by
rank alone.

\subsection{Your fold-created crossings and ours are the same crossings}\label{ss:universal}
One structural observation that I think connects your Remark~4.1 and Remark~6.2 to the
pretzel data. In your setting, the composed curve acquires self-intersections from the fold
singularities of the projection, non-transverse ones that need local Hamiltonian
perturbation. In the pillowcase, the Conway-sum curve acquires its cochain-relevant
self-intersections from resolving the seam fiber circles --- Smith's diagonal cut-and-paste,
one crossing per circle, which is precisely a local perturbation of a non-transverse
composition artifact. All three computed cochains are supported on exactly these
composition-created crossings. Your torus example \cite[\S5]{Z1} reads the same way to me:
the cochain generator sits on the composed curve, born of the fold, and its triangle terms
cancel the composition-side bigon (the deformed differential $\mu_1^{\bc}$ vanishes) ---
cancellation in precisely the sense of
the deficit members of Part~I, and, as far as I can now tell, the first such example in the
literature (a point my accompanying letter returns to, since it obliges a wording change and
a citation in \cite{P2}). So the pretzel family reads as data for your Remark~6.2's
question (a universal, functorial rule attaching cochains): the rule must select
composition-created crossings; at deficit members it is forced (unique minimal cochain), at
excess members rank-matching leaves a $55$-fold ambiguity --- fifty of which share one
common target generator --- so whatever pins the cochain there must come from the naturality
data itself, not from the homology. If your closed-surface pictures show the same
concentration of the cochain on composition-created crossings, that would be worth knowing;
it suggests the universal rule is a statement about the composition map, not about the
curves.

\subsection{Test cases}\label{ss:tests}
The smallest polygons any $k\ge2$ identification must handle, with explicit coordinates and
machine-checkable data (windings, corner sectors, perturbation stability) available on
request or from the repository:
\begin{itemize}
\item the $q=7$ triangle: one $\bc$-vertex at $(\gamma,\theta)\approx(0.05,5.41)$, generator pair
coinciding with a bigon of $\partial_0$; the minimal cancelling configuration;
\item the $q=5$ quadrilateral: $\bc$-vertices $s_A$ at $(\gamma,\theta)\approx(0.028,1.272)$
and $s_B$ at $(3.057,4.981)$, generator vertices $x_4,y_6$; minimal configuration where no
triangle exists --- any theory in which triangles exhaust the $k=2$ corrections is refuted by
this example;
\item any of the fifty-five $q=11$ creating triangles, as the abundance regime.
\end{itemize}
All corner sectors sit in the smooth stratum with the clearances of \S\ref{ss:orbifold}, so
these are simultaneously test cases for the closed-surface statement and for the orbifold
reduction.

\appendix

\section{Map from steps to modules}\label{app:files}
All modules are pure Python~3 standard library; each is a self-checking battery
(\textsc{pass/fail}), and the arXiv ancillary \texttt{README} maps every published number to
the file producing it.

\begin{center}
\begin{tabular}{ll}
\toprule
Step & Module(s)\\
\midrule
curve reconstruction (representation theory) & \texttt{grounded.py}, \texttt{tangles.py}\\
seam-circle resolution; earring & \texttt{resolve.py}, \texttt{earring.py}\\
generators and bigons (\S\ref{sec:naive}) & \texttt{bigons.py}\\
general polygon predicate (\S\ref{ss:tables}) & \texttt{polygons.py}\\
triangle/quadrilateral tables & \texttt{deform.py}, \texttt{deform\_full.py}\\
Maurer--Cartan data (\S\ref{ss:mc}) & \texttt{maurer\_cartan.py}\\
the search (\S\ref{sec:search}) & \texttt{solve\_b2.py}, \texttt{pretzel\_solve.py}\\
perturbation stability (\S\ref{ss:pert}) & \texttt{pert\_check.py}\\
the $q=5$ result battery (\S\ref{ss:q5}) & \texttt{b2\_result.py}\\
the surgery lemma check (\S\ref{ss:surgery}) & \texttt{surgery\_check.py}\\
Alexander recursion (the target, \S\ref{ss:precondition}) & \texttt{skein\_alexander.py}\\
published-number integrity gates & \texttt{check\_paper2\_numbers.py}\\
\bottomrule
\end{tabular}
\end{center}

\begin{thebibliography}{CHKK22}

\bibitem[P1]{P1} B.~J.~Wuebben, \emph{Traceless $\mathrm{SU}(2)$ characters and $\Z/4$
instanton gradings for two-bridge and $(3,n)$-torus knots}, arXiv:2607.26095.

\bibitem[P2]{P2} B.~J.~Wuebben, \emph{The instanton homology of the $(-2,3,q)$ pretzel knots
and computed bounding cochains in the pillowcase}, arXiv:2607.26096.

\bibitem[CHKK22]{CHKK} G.~Cazassus, C.~Herald, P.~Kirk, A.~Kotelskiy, \emph{The
correspondence induced on the pillowcase by the earring tangle}, J.\ Topol.\ \textbf{15}
(2022), 2472--2543. arXiv:2010.04320.

\bibitem[Smi24]{Smi} K.~Smith, \emph{Perturbed traceless $\mathrm{SU}(2)$ character varieties
of tangle sums}, arXiv:2412.06066.

\bibitem[MWW18]{MWW} S.~Ma'u, K.~Wehrheim, C.~Woodward, \emph{$A_\infty$ functors for
Lagrangian correspondences}, Selecta Math.\ \textbf{24} (2018), 1913--2002. arXiv:1601.04919.

\bibitem[PW19]{PW} J.~Palmer, C.~Woodward, \emph{Invariance of immersed Floer cohomology under
Lagrangian surgery}, arXiv:1903.01943.

\bibitem[Zha24]{Z1} Z.~Zhang, \emph{Construction of holomorphic quilts in Cartesian product of
closed surfaces}, arXiv:2409.19744.

\bibitem[Zha25]{Z2} Z.~Zhang, \emph{Uniqueness of holomorphic quilts lifted from holomorphic
bigons on surfaces}, arXiv:2504.09284.

\end{thebibliography}

\bigskip
\noindent{\sc Bernd J. Wuebben, New York, NY,} \texttt{wuebben@gmail.com}

\end{document}
